% ============================================================
% SEÇÃO 5 – DOCUMENTAÇÃO DE ITERAÇÃO
% ============================================================

\section{DOCUMENTAÇÃO DE ITERAÇÃO}

\subsection{Histórias do usuário — Semestre 1}

\subsubsection{Usuário busca evento}

\begin{longtable}{|p{2cm}|p{9cm}|p{2.5cm}|}
  \hline
  \textbf{ID} & \textbf{Descrição} & \textbf{Situação} \\
  \hline
  \endfirsthead
  \hline
  \textbf{ID} & \textbf{Descrição} & \textbf{Situação} \\
  \hline
  \endhead
  USB01 & Como usuário, quero acessar o aplicativo, para visualizar o mapa e buscar eventos próximos. & Completo \\
  \hline
  USB02 & Como usuário, quero buscar eventos próximos à minha localização, para descobrir o que está acontecendo na minha região. & Completo \\
  \hline
  USB03 & Como usuário, quero visualizar os eventos exibidos no mapa, para escolher qual deles me interessa. & Completo \\
  \hline
  USB04 & Como usuário, quero selecionar um evento no mapa, para ver suas informações detalhadas. & Completo \\
  \hline
  USB05 & Como usuário, quero visualizar informações completas sobre o evento (nome, descrição, horário, endereço e categoria), para decidir se quero participar. & Completo \\
  \hline
  USB06 & Como usuário, quero visualizar fotos ou imagens relacionadas ao evento, para entender melhor a atividade antes de participar. & Completo \\
  \hline
\end{longtable}

\subsubsection{Relato de problema no bairro}

\begin{longtable}{|p{2cm}|p{9cm}|p{2.5cm}|}
  \hline
  \textbf{ID} & \textbf{Descrição} & \textbf{Situação} \\
  \hline
  \endfirsthead
  \hline
  \textbf{ID} & \textbf{Descrição} & \textbf{Situação} \\
  \hline
  \endhead
  RPB01 & Como usuário, quero relatar um problema no meu bairro, para que as autoridades responsáveis possam tomar providências. & Completo \\
  \hline
  RPB02 & Como usuário, quero enviar a localização do problema (automaticamente via GPS ou manualmente), para que o órgão responsável saiba exatamente onde agir. & Completo \\
  \hline
  RPB03 & Como usuário, quero adicionar uma descrição explicando o problema encontrado, para que ele seja compreendido corretamente. & Completo \\
  \hline
  RPB04 & Como usuário, quero anexar uma foto do problema, para facilitar a verificação e priorização da solução. & Completo \\
  \hline
  RPB05 & Como usuário, quero receber uma confirmação de que meu relato foi enviado, para ter certeza de que o caso será analisado. & Completo \\
  \hline
  RPB06 & Como usuário, quero acompanhar o status do meu relato, para saber se o problema está sendo resolvido. & Completo \\
  \hline
\end{longtable}

% ──────────────────────────────────────────────────────────────
\subsection{Histórias do usuário — Semestre 2}
% ──────────────────────────────────────────────────────────────

O segundo semestre de desenvolvimento foi focado na consolidação do backend
próprio com autenticação JWT, no redesign visual completo do aplicativo e na
adição de novas funcionalidades de usabilidade e configuração.

\subsubsection{Autenticação com backend próprio}

\begin{longtable}{|p{2cm}|p{9cm}|p{2.5cm}|}
  \hline
  \textbf{ID} & \textbf{Descrição} & \textbf{Situação} \\
  \hline
  \endfirsthead
  \hline
  \textbf{ID} & \textbf{Descrição} & \textbf{Situação} \\
  \hline
  \endhead
  AUTH01 & Como usuário, quero que minha sessão seja mantida automaticamente após o login, sem precisar entrar novamente ao reabrir o aplicativo. & Completo \\
  \hline
  AUTH02 & Como usuário, quero que minha sessão seja renovada automaticamente em segundo plano quando o token expirar, para não ser interrompido durante o uso. & Completo \\
  \hline
  AUTH03 & Como usuário, quero poder alterar minha senha pelo aplicativo, para manter minha conta segura. & Completo \\
  \hline
  AUTH04 & Como usuário, quero poder alterar meu número de telefone pelo aplicativo, para manter meus dados atualizados. & Completo \\
  \hline
\end{longtable}

\subsubsection{Redesign visual e experiência do usuário}

\begin{longtable}{|p{2cm}|p{9cm}|p{2.5cm}|}
  \hline
  \textbf{ID} & \textbf{Descrição} & \textbf{Situação} \\
  \hline
  \endfirsthead
  \hline
  \textbf{ID} & \textbf{Descrição} & \textbf{Situação} \\
  \hline
  \endhead
  UX01 & Como usuário, quero uma interface com modo escuro para usar o aplicativo de forma mais confortável à noite. & Completo \\
  \hline
  UX02 & Como usuário, quero poder alternar entre modo claro e modo escuro nas configurações do aplicativo. & Completo \\
  \hline
  UX03 & Como usuário, quero ver o mascote do LifeCity na tela de splash e onboarding, para uma experiência de entrada mais agradável. & Completo \\
  \hline
  UX04 & Como usuário, quero uma tela de configurações organizada, com opções de aparência, conta e notificações separadas por seções. & Completo \\
  \hline
\end{longtable}

\subsubsection{Gerenciamento de perfil e conta}

\begin{longtable}{|p{2cm}|p{9cm}|p{2.5cm}|}
  \hline
  \textbf{ID} & \textbf{Descrição} & \textbf{Situação} \\
  \hline
  \endfirsthead
  \hline
  \textbf{ID} & \textbf{Descrição} & \textbf{Situação} \\
  \hline
  \endhead
  PROF01 & Como usuário, quero editar meu perfil (nome, foto, CPF, data de nascimento) diretamente pelo aplicativo. & Completo \\
  \hline
  PROF02 & Como usuário, quero visualizar meus eventos e reclamações cadastrados em uma tela separada (``Meus itens''). & Completo \\
  \hline
  PROF03 & Como usuário, quero poder excluir minhas próprias reclamações e eventos. & Completo \\
  \hline
\end{longtable}

% ──────────────────────────────────────────────────────────────
\subsection{Histórias do usuário — Semestre 3}
% ──────────────────────────────────────────────────────────────

O terceiro semestre de desenvolvimento foi focado na gamificação e no engajamento
social, implementando conquistas, notificações, missões automáticas diárias/semanais
e o sistema de equipes colaborativas.

\subsubsection{Gamificação — XP, Níveis e Conquistas}

\begin{longtable}{|p{2cm}|p{9cm}|p{2.5cm}|}
  \hline
  \textbf{ID} & \textbf{Descrição} & \textbf{Situação} \\
  \hline
  \endfirsthead
  \hline
  \textbf{ID} & \textbf{Descrição} & \textbf{Situação} \\
  \hline
  \endhead
  GAM01 & Como usuário, quero acumular XP ao criar reclamações, receber curtidas e comentários, para sentir que minha participação tem valor. & Completo \\
  \hline
  GAM02 & Como usuário, quero subir de nível conforme acumulo XP, para visualizar meu progresso cívico no perfil. & Completo \\
  \hline
  GAM03 & Como usuário, quero desbloquear conquistas automaticamente ao atingir marcos (ex.: criar a primeira reclamação, receber 10 curtidas), para ser recompensado por minhas contribuições. & Completo \\
  \hline
  GAM04 & Como usuário, quero escolher até 3 conquistas para exibir em destaque no meu perfil público, para mostrar minhas realizações. & Completo \\
  \hline
\end{longtable}

\subsubsection{Notificações sociais}

\begin{longtable}{|p{2cm}|p{9cm}|p{2.5cm}|}
  \hline
  \textbf{ID} & \textbf{Descrição} & \textbf{Situação} \\
  \hline
  \endfirsthead
  \hline
  \textbf{ID} & \textbf{Descrição} & \textbf{Situação} \\
  \hline
  \endhead
  NOT01 & Como usuário, quero receber notificações quando alguém curtir ou comentar minhas reclamações, para acompanhar o engajamento. & Completo \\
  \hline
  NOT02 & Como usuário, quero receber notificação ao desbloquear uma conquista, para ser informado da recompensa. & Completo \\
  \hline
  NOT03 & Como usuário, quero receber notificação ao concluir uma missão, para saber que ganhei XP. & Completo \\
  \hline
  NOT04 & Como usuário, quero receber notificação de convite para equipe, para aceitar ou recusar a participação. & Completo \\
  \hline
\end{longtable}

\subsubsection{Missões diárias e semanais}

\begin{longtable}{|p{2cm}|p{9cm}|p{2.5cm}|}
  \hline
  \textbf{ID} & \textbf{Descrição} & \textbf{Situação} \\
  \hline
  \endfirsthead
  \hline
  \textbf{ID} & \textbf{Descrição} & \textbf{Situação} \\
  \hline
  \endhead
  MIS01 & Como usuário, quero receber uma missão diária automaticamente todo dia, para ter um objetivo simples de engajamento. & Completo \\
  \hline
  MIS02 & Como usuário, quero receber uma missão semanal automaticamente toda semana, para ter um desafio maior a cumprir. & Completo \\
  \hline
  MIS03 & Como usuário, quero visualizar o progresso das minhas missões em tempo real com uma barra de progresso, para saber o quanto falta para concluir. & Completo \\
  \hline
  MIS04 & Como usuário, quero ganhar XP ao concluir uma missão, para ser recompensado pelo meu engajamento. & Completo \\
  \hline
\end{longtable}

\subsubsection{Equipes colaborativas}

\begin{longtable}{|p{2cm}|p{9cm}|p{2.5cm}|}
  \hline
  \textbf{ID} & \textbf{Descrição} & \textbf{Situação} \\
  \hline
  \endfirsthead
  \hline
  \textbf{ID} & \textbf{Descrição} & \textbf{Situação} \\
  \hline
  \endhead
  EQ01 & Como usuário, quero criar uma equipe com nome próprio e convidar amigos, para colaborar em missões semanais. & Completo \\
  \hline
  EQ02 & Como usuário, quero aceitar ou recusar convites para entrar em equipes, para controlar minha participação. & Completo \\
  \hline
  EQ03 & Como usuário, quero visualizar os membros da minha equipe, o XP coletivo e o status de cada integrante, para acompanhar o desempenho do grupo. & Completo \\
  \hline
  EQ04 & Como usuário, quero receber bônus de XP semanal proporcional ao desempenho dos meus colegas de equipe, para ser recompensado pela colaboração. & Completo \\
  \hline
\end{longtable}
